\section{演習の操作を別の表現で見直す}\label{節：演習の操作を別の表現で見直す}
  新しい道具を学んだ後に,すでに解いた問題へ戻ろう.
  ここでは\bookproblemref{practice-entrance}{3}{発展問3},
  \bookproblemref{practice-entrance}{4}{問4},
  \bookproblemref{practice-entrance}{9}{問9}で行った操作を,
  \bookterm{generating}{母関数}・回転・\bookterm{invariant}{不変量}という見方で確かめる.

  \subsection{畳み込み和が一定になる理由}
  発展問3の$c_0=1$, $c_{n+1}=\frac{2n+1}{2n+2}c_n$から,
  \[
    c_n=\frac1{4^n}\binom{2n}{n}
  \]
  を得た.その\bookterm{generating}{母関数}を$C(x)=\sum_{n\ge0}c_nx^n$とする.
  まず,級数の各項を微分する形式的な操作を使う.
  収束する関数の微分を仮定するのではなく,
  $x^n$を$nx^{n-1}$へ移して係数を並べ直す操作として読む.
  元の\bookterm{recurrence}{漸化式}を$(n+1)c_{n+1}=(n+\frac12)c_n$と書けば,
  \[
    (1-x)C'(x)=\frac12C(x)
  \]
  を得る.積の係数は有限和なので,係数を比較することで通常の積の微分法則も使える.
  $D(x)=C(x)^2=\sum_{n\ge0}d_nx^n$と置くと,
  \[
    (1-x)D'(x)=2C(x)(1-x)C'(x)=D(x).
  \]
  $x^n$の係数を比較すれば$(n+1)d_{n+1}-nd_n=d_n$,すなわち$d_{n+1}=d_n$である.
  $d_0=c_0^2=1$なので,
  \[
    C(x)^2=\frac1{1-x}=\sum_{n\ge0}x^n
  \]
  の$x^n$の係数から
  \[
    \sum_{k=0}^n c_kc_{n-k}=1
  \]
  を得る.問題の解法は有限和を変形し,この解法は\bookterm{sequence}{数列}全体の積へ移した.
  どちらも同じ\bookterm{convolution}{畳み込み}を扱っている.

  \subsection{中点の更新を,回転と面積から見直す}
  発展問4では,平行四辺形の各辺上で同じ比率$t$ $(0<t<1)$の分点を取り,
  新しい平行四辺形を作った.連続する辺ベクトルを$u,v$とする.
  頂点の更新を引き算すれば,
  \[
    u'=(1-t)u+tv,\qquad v'=-tu+(1-t)v
  \]
  である.従って辺の組を混ぜる\bookterm{matrix}{行列}は
  \[
    M=\begin{pmatrix}1-t&t\\-t&1-t\end{pmatrix}
  \]
  となる.\bookterm{matrix}{行列}式は$\rho=(1-t)^2+t^2$だから,
  平行四辺形の面積は一回ごとに$\rho$倍となる.
  \bookterm{eigenvalue}{固有値}は$1-t\pm it$であり,その絶対値の二乗も$\rho$である.

  正方形で,向きを反時計回りに取った辺を複素数として$v=iu$と書けるときは,
  \[
    u'=((1-t)+it)u
  \]
  である.$t=\frac12$なら倍率は$\frac1{\sqrt2}$,回転角は$\frac\pi4$になる.
  これが問題の中点更新の辺長と向きに一致する.
  \begin{bookpoint}
    \textbf{\bookterm{matrix}{行列}が混ぜる対象を確認する}\par
    一般の平行四辺形では$u,v$は直交する等長の辺ではない.
    辺の組を混ぜる\bookterm{matrix}{行列}を,そのまま平面上の相似変換と読むことはできない.
    面積の倍率の説明は一般の場合に使え,上の回転の説明には正方形という条件を使う.
  \end{bookpoint}

  \subsection{整数方程式を保つ更新を,不変量から見る}
  発展問9では,$u_0=3$, $u_1=3$, $u_{n+2}=3u_{n+1}-u_n$から,
  $(3,u_n,u_{n+1})$が$x^2+y^2+z^2=xyz$の解になることを示した.
  初期値では
  \[
    u_n^2+u_{n+1}^2+9=3u_nu_{n+1}
  \]
  が成り立つ.一般の$x,y$に対して
  \[
    y^2+(3y-x)^2-3y(3y-x)=x^2+y^2-3xy
  \]
  だから,左辺の差は更新の前後で変わらず$-9$に保たれる.
  その結果,$u_nu_{n+2}=u_{n+1}^2+9$となり,
  \[
    u_{n+2}=\frac{u_{n+1}^2+9}{u_n}
  \]
  とも書ける.正値性は$u_0=u_1=3$と,
  $u_2=6$以降の増加から確認できる.
  \sectionref*{節：不変量から線形の関係}の定数$1$を$9$へ変えた更新だった.
  元の問題では\bookterm{linear}{線形}の更新から保存を示し,本編の新しい例では保存から\bookterm{linear}{線形}の更新を取り出した.
  同じ関係を二つの向きから使っている.

  \begin{bookpoint}
    \textbf{別解から何を持ち帰るか}\par
    計算が短くなるだけでなく,足す意味・面積の意味・保存される関係が見える.
    新しい表現を使った後は,元の問題の条件と答えへ戻って対応を確かめよう.
  \end{bookpoint}
